\documentclass[11pt,a4paper]{article}

\usepackage[T1]{fontenc}
\usepackage[utf8]{inputenc}
\usepackage[spanish]{babel}
\usepackage[a4paper,margin=2.6cm]{geometry}
\usepackage{xcolor}
\usepackage{titlesec}
\usepackage[hidelinks]{hyperref}
\usepackage{tikz}
\usetikzlibrary{arrows.meta}
\usepackage{listings}

\definecolor{UPCBlue}{HTML}{123B66}
\definecolor{UPCLightBlue}{HTML}{D9E8F5}

\setlength{\parindent}{0pt}
\setlength{\parskip}{0.35em}
\setcounter{secnumdepth}{3}
\setcounter{tocdepth}{3}
\titleformat{\section}{\Large\bfseries\color{UPCBlue}}{\thesection}{0.7em}{}
\titleformat{\subsection}{\large\bfseries\color{UPCBlue}}{\thesubsection}{0.7em}{}
\titleformat{\subsubsection}{\normalsize\bfseries}{\thesubsubsection}{0.7em}{}
\lstdefinestyle{salidaclip}{
    basicstyle=\ttfamily\scriptsize,
    breaklines=true,
    breakatwhitespace=false,
    columns=fullflexible,
    keepspaces=true,
    showstringspaces=false,
    frame=single,
    framerule=0.3pt,
    rulecolor=\color{UPCBlue},
    xleftmargin=0.5em,
    xrightmargin=0.5em
}

\newcommand{\tituloPractica}{Sistema experto de\\[0.18em] recomendaci\'on de viajes}
\newcommand{\subtituloPractica}{Pr\'actica de Sistemas Basados en el Conocimiento}
\newcommand{\asignatura}{Inteligencia Artificial}
\newcommand{\curso}{Curso 2025/2026 -- 2Q}

\begin{document}

\hypersetup{pageanchor=false}
\begin{titlepage}
    \thispagestyle{empty}
    \color{UPCBlue}
    \vspace*{1.2cm}

    {\Large Universitat Polit\`ecnica de Catalunya\par}
    {\large Grau en Enginyeria Inform\`atica\par}
    {\large \asignatura\par}

    \vspace{2.2cm}
    {\color{UPCLightBlue}\rule{\textwidth}{1.2pt}\par}
    \vspace{0.7cm}
    {\Huge\bfseries \tituloPractica\par}
    \vspace{0.45cm}
    {\LARGE \subtituloPractica\par}
    \vspace{0.7cm}
    {\color{UPCLightBlue}\rule{\textwidth}{1.2pt}\par}

    \vspace{2.1cm}
    {\large Documentaci\'on de la pr\'actica\par}
    {\large \curso\par}

    \vfill
    {\large\bfseries Autores\par}
    \vspace{0.35cm}
    {\large Bernat Llorens\par}
    {\large Santiago Fernandez\par}
    {\large Roger Garc\'ia\par}

    \vspace{1.0cm}
\end{titlepage}

\hypersetup{pageanchor=true}
\pagenumbering{roman}
\setcounter{page}{1}
\tableofcontents
\newpage

\pagenumbering{arabic}
\setcounter{page}{1}

\section{Introducci\'on}
\subsection{Contexto de la pr\'actica y motivaci\'on del sistema}
En esta pr\'actica hemos desarrollado un Sistema Basado en el Conocimiento
orientado a recomendar viajes personalizados para la agencia ``Al fin del mundo
y m\'as all\'a''. El problema no consiste solamente en escoger un destino bonito:
un viaje completo depende del tipo de usuario, de sus restricciones y de varias
decisiones que est\'an relacionadas entre s\'i, como el n\'umero de ciudades, los
d\'ias de estancia, el alojamiento, los transportes entre etapas y las visitas
que se proponen en cada ciudad.

Por ejemplo, no tiene sentido tratar igual una escapada de aniversario, un
viaje familiar con ni\~nos o un grupo de amigos que viaja por fin de curso. Dos
usuarios pueden tener el mismo presupuesto y aun as\'i necesitar propuestas muy
distintas. Adem\'as, algunas condiciones se tienen que cumplir de forma
obligatoria, como no superar el presupuesto o evitar el avi\'on si el usuario lo
pide, mientras que otras sirven para comparar alternativas, como priorizar el
tren o preferir ciudades menos masificadas.

Por este motivo hemos planteado el sistema como un SBC implementado en CLIPS.
La idea es representar de forma expl\'icita el conocimiento del dominio y usar
reglas para ir pasando de la informaci\'on que nos da el usuario a dos
recomendaciones completas y justificadas. De esta manera no nos limitamos a
mostrar una lista fija de viajes, sino que el sistema puede preguntar, deducir,
filtrar combinaciones inviables y comparar las soluciones factibles.


\section{Identificaci\'on del problema}
\subsection{Descripci\'on del problema de recomendaci\'on de viajes}
El problema que queremos resolver es la personalizaci\'on de viajes de varias
ciudades a partir de la informaci\'on de un cliente. El sistema tiene que
obtener primero los datos que condicionan la recomendaci\'on: edad del viajero
principal, tipo de grupo, si viajan ni\~nos, si existe un evento especial,
objetivo del viaje, presupuesto, duraci\'on, n\'umero de ciudades, calidad
m\'inima de alojamiento y preferencias sobre transporte o destinos menos
conocidos.

Con esos datos se debe construir una recomendaci\'on concreta. En nuestro caso,
un plan no es solo una ciudad candidata: es un itinerario de dos o tres
ciudades con los d\'ias repartidos entre ellas, los transportes necesarios para
salir del origen, moverse entre etapas y volver, un alojamiento por ciudad,
visitas recomendadas y un precio total estimado por persona.

El problema tiene una parte de an\'alisis del usuario y una parte de
construcci\'on de la soluci\'on. Primero hay que interpretar el perfil y completar
algunas decisiones razonables cuando la entrada lo permite. Por ejemplo, una
boda o un aniversario apuntan a un viaje rom\'antico, y un viaje corto no deber\'ia
acabar en una ruta demasiado extensa. Despu\'es hay que construir planes que
respeten las restricciones y escoger entre las alternativas posibles.

\subsection{Alcance, supuestos y l\'imites del sistema}
Hemos acotado el alcance para poder razonar con suficiente detalle sobre los
planes sin convertir la pr\'actica en un buscador de reservas real. El sistema
trabaja con una base de conocimiento cerrada de ciudades, alojamientos,
actividades, gu\'ias de visita y conexiones de transporte. A partir de esa base
genera viajes de dos o tres ciudades y considera un origen gen\'erico para las
conexiones de ida y vuelta.

Los precios son estimaciones por persona y se usan para comparar y validar
planes dentro del sistema. Para construirlos hemos tomado como referencia la
informaci\'on consultada sobre destinos, alojamientos, visitas y conexiones,
pero no pretendemos que la base de conocimiento reproduzca precios de mercado,
disponibilidad hotelera o rutas comerciales actualizadas. En cambio, s\'i
hemos intentado mantener criterios coherentes: los transportes lejanos son
m\'as caros, la calidad del alojamiento influye en el precio y un viaje corto
limita las combinaciones razonables.

El sistema tampoco intenta resolver todos los detalles de la organizaci\'on real
de un viaje. No gestiona fechas concretas, horarios, visados, meteorolog\'ia,
reservas ni incidencias durante el recorrido. Su objetivo es producir
recomendaciones explicables a partir del conocimiento que hemos representado,
y detectar cuando las restricciones del usuario no dejan ninguna alternativa
factible con la base disponible.

\subsection{Viabilidad de construir un Sistema Basado en el Conocimiento}
Consideramos que el problema es adecuado para construir un SBC porque no basta
con aplicar una f\'ormula sobre los datos del usuario. Para recomendar un viaje
hay que combinar conocimiento del dominio con decisiones de experto: distinguir
restricciones de preferencias, interpretar algunos perfiles, descartar
combinaciones poco razonables y comparar rutas que ya son factibles.

Adem\'as, el problema se puede describir mediante objetos y relaciones claras.
Tenemos usuarios, ciudades, alojamientos, actividades, conexiones y planes.
Sobre estos conceptos aparecen reglas bastante naturales: si se viaja con
ni\~nos se deben evitar recursos no aptos, si se quiere evitar el avi\'on se
descartan rutas que lo usan, si el presupuesto no cubre el precio calculado el
plan no sirve, y si el usuario prefiere tren esa caracter\'istica puede mejorar
la puntuaci\'on de una alternativa sin convertirla en obligatoria.

Esta separaci\'on entre conocimiento representado y reglas de razonamiento nos
permite justificar las recomendaciones y tambi\'en explicar por qu\'e puede no
haber soluci\'on. Por eso un lenguaje de reglas como CLIPS encaja bien con el
problema: el sistema puede ir acumulando hechos sobre el usuario y los planes,
activar reglas cuando se cumplen sus condiciones y controlar el razonamiento
por fases hasta llegar a la salida.

\subsection{Fuentes de conocimiento utilizadas}
Para construir el sistema hemos usado varias fuentes de conocimiento. La
primera ha sido el propio enunciado de la pr\'actica, que fija el problema que
hay que resolver, los datos relevantes que se deben considerar y la diferencia
entre restricciones obligatorias y preferencias que permiten comparar
soluciones.

Tambi\'en hemos usado el material de la asignatura sobre representaci\'on del
conocimiento, ontolog\'ias y arquitectura de los SBC. Ese material nos ha
servido para estructurar el dominio, construir la ontolog\'ia, separar la base de
conocimiento del razonamiento y organizar la documentaci\'on siguiendo las
fases de la ingenier\'ia del conocimiento.

Para la parte concreta de viajes hemos consultado informaci\'on tur\'istica de
cat\'alogos y bases de datos de agencias de viajes, alojamientos, actividades y
transportes. A partir de esa informaci\'on hemos seleccionado destinos,
alojamientos, visitas y conexiones que nos permitieran construir una base de
conocimiento variada, y los hemos adaptado al nivel de detalle que necesita el
sistema con precios estimados, niveles de calidad, popularidad y etiquetas
tur\'isticas. El objetivo de esos datos no es sustituir a una agencia real,
sino representar suficiente conocimiento del dominio para que las reglas
produzcan itinerarios distintos y casos de prueba relevantes.

Por \'ultimo, durante la conceptualizaci\'on hemos usado un modelo de lenguaje
como experto simulado para contrastar criterios de recomendaci\'on. Esa
conversaci\'on nos ha servido como apoyo para decidir qu\'e conocimiento del
dominio era \'util convertir en reglas o preferencias. Documentaremos el prompt,
las preguntas relevantes y el conocimiento que hemos aprovechado en el
apartado espec\'ifico de conceptualizaci\'on.

\subsection{Objetivos del sistema}
Los objetivos que nos hemos marcado para el sistema son los siguientes:
\begin{itemize}
    \item Obtener del usuario la informaci\'on necesaria para personalizar el
    viaje y no basar la recomendaci\'on solo en el presupuesto o en un destino
    aislado.
    \item Deducir parte del perfil del viajero y aplicar conocimiento de sentido
    com\'un cuando las respuestas permiten completar la interpretaci\'on del caso.
    \item Representar de forma expl\'icita las restricciones del viaje y
    descartar cualquier plan que no las cumpla.
    \item Tratar las preferencias como criterios de comparaci\'on para puntuar
    soluciones factibles y escoger alternativas m\'as adecuadas.
    \item Generar viajes multi-ciudad con alojamiento, transportes, visitas,
    duraci\'on y precio total estimado.
    \item Recomendar dos planes diferentes cuando la base de conocimiento lo
    permita, e informar claramente si no existe ninguna recomendaci\'on factible.
\end{itemize}

\subsection{Resultados esperados y requisitos de la recomendaci\'on}
Al final de la ejecuci\'on esperamos que el usuario reciba dos recomendaciones
de viaje suficientemente concretas para entender qu\'e propone el sistema y por
qu\'e. Cada plan debe indicar:
\begin{itemize}
    \item el precio total estimado por persona;
    \item la duraci\'on total del viaje;
    \item las ciudades visitadas y los d\'ias asignados a cada una;
    \item las visitas recomendadas en cada ciudad;
    \item el alojamiento escogido en cada etapa;
    \item los transportes desde el origen, entre ciudades y de vuelta;
    \item las preferencias del usuario que se han cumplido.
\end{itemize}

Las dos recomendaciones tienen que ser alternativas reales y no dos
impresiones del mismo plan. Por eso se busca que visiten ciudades diferentes y
se seleccionan sobre el conjunto de candidatos que ya cumplen presupuesto,
duraci\'on, calidad m\'inima y restricciones de transporte. Si las condiciones
del usuario son incompatibles con la base de conocimiento disponible, el
resultado esperado pasa a ser un mensaje expl\'icito de no soluci\'on en lugar de
forzar una recomendaci\'on incorrecta.

\section{Conceptualizaci\'on}
\subsection{Descripci\'on del dominio de viajes personalizados}
En la fase de conceptualizaci\'on hemos pasado del enunciado general de
``recomendar viajes'' a una descripci\'on del dominio que el sistema pueda
tratar. Para nosotros un viaje personalizado es una propuesta que relaciona
dos partes: por un lado el cliente, con sus condiciones y sus gustos; por el
otro los recursos tur\'isticos disponibles, que son destinos, alojamientos,
visitas y transportes.

La recomendaci\'on no se decide con un \'unico dato. Un presupuesto alto no
convierte cualquier ruta en buena, igual que escoger un motivo cultural no
anula la duraci\'on, el tipo de grupo o la calidad m\'inima de alojamiento. Por
eso el dominio mezcla restricciones que se deben cumplir con criterios que
permiten escoger la alternativa m\'as adecuada entre varias que ya son
factibles.

El resultado del dominio tampoco es un destino aislado. La soluci\'on que
queremos construir es un itinerario de dos o tres ciudades con etapas
ordenadas. Cada etapa necesita una ciudad, unos d\'ias de estancia, un
alojamiento y unas visitas concretas; adem\'as el plan completo necesita los
transportes de ida, entre ciudades y de vuelta, un precio estimado y una
explicaci\'on breve de las preferencias que satisface.

Con esta forma de ver el problema, el conocimiento que necesitamos no se
limita a almacenar datos de turismo. Tambi\'en hace falta representar
decisiones de sentido com\'un que usar\'ia una agencia al hablar con un cliente:
qu\'e destinos encajan mejor con un objetivo rom\'antico o de descanso, qu\'e
cambia cuando viajan ni\~nos, cu\'ando una ruta lejana deja de tener sentido
para pocos d\'ias y c\'omo comparar dos viajes que cumplen las condiciones
obligatorias.

\subsection{Conceptos principales del dominio}
\subsubsection{Usuario, perfil de viaje y grupo viajero}
El punto de partida es el usuario. No lo tratamos solo como la persona
que ejecuta el programa, sino como el viajero o grupo para el que se prepara la
recomendaci\'on. De \'el nos interesan la edad del viajero principal, la
composici\'on del grupo, si viajan ni\~nos y si el viaje est\'a ligado a un evento
especial como una boda, un aniversario o un fin de curso.

Sobre el usuario tambi\'en aparecen las condiciones propias del viaje:
presupuesto m\'aximo, duraci\'on disponible, n\'umero de ciudades que se quiere
visitar, calidad m\'inima del alojamiento y decisiones sobre el transporte. A
eso se a\~naden preferencias m\'as blandas, por ejemplo querer usar tren cuando
sea razonable o preferir ciudades menos masificadas.

El perfil de viaje es la lectura que hacemos de esa informaci\'on. El
usuario puede indicar directamente su motivo principal, pero hay casos donde
el sistema puede completar la interpretaci\'on: una boda o un aniversario
apuntan a un viaje rom\'antico, un fin de curso a uno de diversi\'on y un grupo
con ni\~nos obliga a pensar en destinos y alojamientos adecuados. Esta
separaci\'on entre datos preguntados y perfil interpretado es importante porque
es donde aparece parte del conocimiento experto.

\subsubsection{Ciudades y caracter\'isticas tur\'isticas}
Las ciudades son los destinos que pueden formar parte de una ruta.
Para recomendarlas no basta con conocer su nombre. Necesitamos saber su pa\'is,
la zona o continente, un coste diario estimado, su nivel de seguridad, si son
aptas para viajar con ni\~nos, su popularidad y un conjunto de etiquetas
tur\'isticas.

Las etiquetas resumen qu\'e puede aportar cada destino: cultural, rom\'antico,
descanso, diversi\'on, naturaleza, playa, gastronom\'ia o un car\'acter m\'as
ic\'onico. Gracias a ellas se puede razonar sobre la compatibilidad entre el
objetivo del viaje y la ciudad sin convertir el sistema en una lista cerrada
de respuestas prefijadas para cada caso.

\subsubsection{Alojamientos y calidad de estancia}
Un alojamiento es un recurso asociado a una ciudad. En el dominio nos
interesan su nombre, categor\'ia, calidad, precio por noche y compatibilidad con
ni\~nos. Estos datos permiten cubrir dos necesidades diferentes: comprobar que
se respeta la calidad m\'inima que exige el usuario y construir el precio total
del plan con un coste de estancia concreto.

La calidad del alojamiento se trata con cuidado porque puede actuar como
restricci\'on o como margen de ajuste. Si el usuario pide una calidad m\'inima,
el viaje no deber\'ia bajar de ella sin permiso. En cambio, cuando declara que
acepta sacrificar calidad para ajustar presupuesto, se puede permitir una
flexibilizaci\'on controlada en vez de inventar una soluci\'on que no existe.

\subsubsection{Actividades, visitas y gu\'ias de visita}
Las actividades representan visitas concretas que se pueden hacer en
una ciudad. Cada una tiene nombre, tipo, duraci\'on, coste, prioridad y
compatibilidad familiar. Nos sirven para que la base de conocimiento no se
quede en frases gen\'ericas del estilo ``hacer turismo por la ciudad''.

Para la recomendaci\'on final usamos adem\'as gu\'ias de visita. Una
gu\'ia agrupa visitas adecuadas para una ciudad y un motivo de viaje, indica el
m\'inimo de d\'ias que necesita y aporta un coste total de esas visitas. De esta
manera, la etapa de una ciudad ya lleva una propuesta concreta distinta si el
usuario viaja por cultura, descanso, naturaleza, diversi\'on o romanticismo.

\subsubsection{Conexiones, transportes y etapas del recorrido}
Las conexiones expresan qu\'e transportes est\'an disponibles entre el
origen y una ciudad o entre dos ciudades de la ruta. De cada conexi\'on
necesitamos el medio, el coste, la duraci\'on aproximada y el \'ambito del tramo.
El \'ambito permite distinguir, por ejemplo, un salto intercontinental de un
desplazamiento local o europeo.

Una etapa es la parte del viaje que ocurre en una ciudad. Aunque en la
implementaci\'on final la etapa queda embebida dentro del plan candidato, en la
conceptualizaci\'on nos ayuda a ordenar la soluci\'on: una ruta est\'a formada por
etapas y conexiones entre ellas.

\subsubsection{Restricciones obligatorias y preferencias comparativas}
Una decisi\'on central de la conceptualizaci\'on ha sido separar
restricciones de preferencias. Las restricciones son las
condiciones que no queremos violar al generar una recomendaci\'on: no superar
el presupuesto, respetar la duraci\'on disponible, aceptar solo el n\'umero de
ciudades permitido, evitar el avi\'on cuando se pide, mantener la calidad
m\'inima y exigir compatibilidad familiar cuando viajan ni\~nos.

Las preferencias no tienen el mismo papel. Se usan para comparar soluciones
que ya han pasado los filtros obligatorios. Por ejemplo, una ruta puede recibir
mejor valoraci\'on si usa tren cuando el usuario lo prefiere, si incluye
ciudades menos conocidas, si mejora la calidad m\'inima pedida o si deja margen
de presupuesto. As\'i evitamos descartar un viaje v\'alido solo porque no cumple
un gusto opcional.

\subsubsection{Planes candidatos y planes recomendados}
Un plan candidato es una combinaci\'on que el sistema ha conseguido
construir con ciudades, d\'ias, alojamientos, transportes, visitas, precio y
puntuaci\'on. Solo llega a ser candidato cuando respeta las restricciones que
se han comprobado durante la generaci\'on.

Un plan recomendado es uno de los candidatos que se escoge al final.
El sistema no debe quedarse con el primero que aparezca, sino comparar las
alternativas disponibles y elegir dos rutas distintas siempre que la base de
conocimiento lo permita. Esta diferencia entre candidato y recomendaci\'on nos
permite separar la construcci\'on de soluciones de la elecci\'on final.

\subsection{Conocimiento experto extra\'ido del dominio}
\subsubsection{Reglas de sentido com\'un sobre perfiles de usuario}
Durante la conceptualizaci\'on hemos identificado reglas sencillas que no son
datos tur\'isticos puros, sino decisiones que usar\'ia una persona que prepara
viajes. Algunos ejemplos que hemos incorporado son:
\begin{itemize}
    \item Una boda o un aniversario se interpreta como una situaci\'on donde
    tiene sentido priorizar un perfil rom\'antico.
    \item Un fin de curso encaja mejor con un perfil de ocio y diversi\'on.
    \item Si viajan ni\~nos, hay que dar importancia a destinos y alojamientos
    aptos y evitar propuestas que dependan de recursos no familiares.
    \item Si un usuario que viaja solo no declara un motivo claro, una ruta
    cultural es un punto de partida razonable y explicable.
\end{itemize}

Estas reglas no sustituyen las respuestas del usuario cuando ya son claras.
Sirven para completar informaci\'on poco concreta y para que el sistema pueda
seguir razonando con un perfil de viaje definido.

\subsubsection{Criterios de compatibilidad de destinos, alojamientos y transportes}
El siguiente bloque de conocimiento experto consiste en decidir qu\'e recursos
pueden convivir dentro de una misma recomendaci\'on. En destinos se comprueba
que el motivo del viaje encaje con las etiquetas de las ciudades. En
alojamientos se comprueba la calidad solicitada y, cuando hay ni\~nos, que el
recurso sea apto. En las visitas se escoge una gu\'ia coherente con la ciudad,
el objetivo del viaje y los d\'ias asignados.

Con los transportes aplicamos criterios parecidos. Si el usuario quiere evitar
avi\'on, una ruta que lo necesita deja de ser compatible. Adem\'as, una ruta muy
lejana no es buena candidata para una escapada demasiado corta: aunque exista
una conexi\'on, el tiempo disponible cambia lo que es razonable recomendar.

\subsubsection{Criterios para comparar soluciones factibles}
Una vez un plan cumple las condiciones obligatorias, todav\'ia puede ser mejor
o peor que otro. Para compararlos hemos identificado varios criterios:
\begin{itemize}
    \item La adecuaci\'on base de las ciudades que forman la ruta.
    \item El uso del tren cuando el usuario lo ha marcado como preferencia.
    \item La presencia de ciudades menos conocidas o menos masificadas si se
    han pedido.
    \item La calidad extra del alojamiento por encima del m\'inimo solicitado.
    \item El margen que queda respecto al presupuesto m\'aximo.
\end{itemize}

Este conocimiento se traduce despu\'es en una puntuaci\'on heur\'istica, pero en
esta fase lo importante es la idea: primero se garantizan las restricciones y
despu\'es se ordenan las alternativas usando preferencias.

\subsubsection{Tratamiento de restricciones incompatibles y no soluci\'on}
Tambi\'en hemos considerado el caso en el que el cliente pide una combinaci\'on
que la base de conocimiento no puede satisfacer. Por ejemplo, un presupuesto
muy bajo junto con una calidad alta, pocos d\'ias y transportes muy limitados
puede dejar el conjunto de candidatos vac\'io.

En ese caso el comportamiento correcto no es relajar condiciones a escondidas
ni imprimir un viaje que no cumple lo pedido. La conclusi\'on del razonamiento
debe ser que no se ha encontrado una recomendaci\'on factible y conviene
indicar qu\'e tipo de restricci\'on podr\'ia revisarse.

\subsection{Conceptualizaci\'on apoyada por un modelo de lenguaje}
\subsubsection{Contexto y prompt de elicitaci\'on de conocimiento}
Como apoyo a esta fase hemos usado un modelo de lenguaje en el papel de
experto simulado en recomendaci\'on de viajes. No le pedimos que escribiera la
ontolog\'ia ni las reglas de CLIPS, sino que respondiera como lo har\'ia un
agente de viajes al explicar sus criterios a un ingeniero del conocimiento.
El contexto base utilizado fue el siguiente:

\begin{quote}
\small\itshape You are going to help a knowledge engineer perform knowledge
elicitation for building a rule-based expert system for touristic travel
recommendation. You are an expert travel agent. Explain the criteria you use
to match traveller profiles, constraints and preferences with destinations,
accommodation, transport and activities. Answers must be direct and useful for
building rules.
\end{quote}

Con este contexto busc\'abamos obtener conocimiento del dominio y del proceso de
decisi\'on, no una soluci\'on ya programada. Despu\'es contrastamos las ideas
\'utiles con el alcance del enunciado y con lo que realmente quer\'iamos
representar en nuestro sistema.

\subsubsection{Preguntas realizadas al experto simulado}
Las preguntas m\'as relevantes que usamos para orientar la extracci\'on de
conocimiento fueron:
\begin{itemize}
    \item \textit{When travelling with children, what kind of accommodation
    and destinations are safer?}
    \item \textit{When the budget is low, is it better to reduce days, quality
    or transportation cost?}
    \item \textit{For a romantic trip, which city features matter most?}
    \item \textit{When should a short trip avoid far-away cities?}
    \item \textit{How do you compare two valid routes when both satisfy hard
    constraints?}
\end{itemize}

Las preguntas est\'an planteadas para sacar criterios de decisi\'on: qu\'e datos
del usuario importan, qu\'e condiciones son obligatorias, cu\'ando se puede
flexibilizar algo y qu\'e elementos sirven para comparar rutas posibles.

\subsubsection{Conocimiento incorporado, adaptado o descartado}
De esa elicitaci\'on nos quedamos con varias ideas que encajaban con la
pr\'actica. Incorporamos la importancia de la compatibilidad familiar cuando
viajan ni\~nos, la lectura rom\'antica de viajes de pareja ligados a boda o
aniversario, la limitaci\'on de rutas lejanas en viajes cortos y la separaci\'on
entre condiciones duras y preferencias.

Otras ideas se adaptaron al tama\~no del sistema. Por ejemplo, el experto puede
hablar de ajustar muchos gastos distintos, pero nosotros hemos representado el
presupuesto con precios estimados de transporte, estancia, coste diario y
visitas, y solo flexibilizamos la calidad del alojamiento cuando el usuario lo
permite. Tambi\'en convertimos la comparaci\'on informal entre rutas en criterios
concretos que luego pueden puntuar planes factibles.

Por \'ultimo, dejamos fuera detalles que podr\'ian ser realistas en una agencia
pero no forman parte de esta base de conocimiento: disponibilidad de reservas,
fechas exactas, climatolog\'ia, documentaci\'on de entrada a pa\'ises o cambios de
precio en tiempo real. Incluirlos habr\'ia ampliado el dominio sin mejorar la
resoluci\'on del problema que se pide en la pr\'actica.

\subsection{Descomposici\'on del problema en subproblemas}
\subsubsection{Adquisici\'on de datos del usuario}
El primer subproblema es recoger la informaci\'on necesaria del cliente. Se
pregunta por su grupo, ni\~nos, evento, motivo, presupuesto, duraci\'on, n\'umero
de ciudades, preferencias de transporte, calidad m\'inima, posibles sacrificios
para ajustar coste, inter\'es por ciudades menos conocidas y ritmo del viaje.
Sin esta parte no hay un caso concreto sobre el que razonar.

\subsubsection{Deducci\'on del perfil, restricciones y preferencias}
Con los datos anteriores se completa el perfil de viaje y se organizan las
decisiones que afectar\'an al resto del proceso. Aqu\'i se deduce el motivo si
falta o si un evento lo aclara, se traduce el presupuesto a un nivel orientativo
y se fijan las restricciones y preferencias que tendr\'an un papel distinto en
la generaci\'on.

\subsubsection{Filtrado de recursos compatibles}
Antes de construir planes completos hay que evitar combinaciones que ya se
sabe que no encajan. Se filtran ciudades por motivo y compatibilidad familiar,
alojamientos por calidad y aptitud para ni\~nos, gu\'ias por motivo y d\'ias
disponibles, y conexiones por las restricciones de transporte y distancia.

\subsubsection{Generaci\'on de itinerarios candidatos}
Con los recursos compatibles se proponen itinerarios de dos o tres ciudades.
Cada candidato reparte los d\'ias de estancia entre sus etapas, incorpora los
transportes necesarios, selecciona alojamientos y arrastra las visitas que se
mostrar\'an al final. Este subproblema construye el espacio de soluciones sobre
el que luego se decidir\'a.

\subsubsection{C\'alculo de precio y puntuaci\'on de alternativas}
Cada itinerario necesita un precio total estimado por persona. Este precio se
forma con transportes, coste de alojamientos, coste diario de las ciudades y
coste de visitas. Si queda por encima del presupuesto, el candidato no es
v\'alido. Si lo respeta, se calcula una puntuaci\'on para reflejar su adecuaci\'on
y las preferencias opcionales que satisface.

\subsubsection{Selecci\'on de dos planes con ciudades diferentes}
Una vez generados los candidatos se escoge el mejor valorado y se busca una
segunda alternativa que no sea una copia con otro alojamiento o con una
variaci\'on irrelevante. La prioridad es encontrar una ruta sin solape de
ciudades; si la base no permite tanto, al menos debe ser una ruta distinta.

\subsubsection{Explicaci\'on del resultado o del caso sin recomendaci\'on}
El \'ultimo subproblema es comunicar la conclusi\'on. Cuando hay planes, se
imprimen de forma completa con ciudades, d\'ias, visitas, alojamientos,
transportes, precio, puntuaci\'on y preferencias cumplidas. Cuando no aparece
ning\'un candidato factible, se explica que las restricciones no se pueden
satisfacer con la base de conocimiento actual en vez de ocultar el fallo.

\subsection{Descripci\'on informal del proceso de resoluci\'on}
Visto de principio a fin, el sistema empieza hablando con el usuario para
convertir una petici\'on abierta en un caso concreto. Despu\'es interpreta ese
caso y completa decisiones sencillas de perfil. Con esa informaci\'on recorre la
base de conocimiento para proponer rutas que tengan ciudades compatibles,
transportes posibles, alojamientos aceptables y visitas coherentes.

Cada ruta propuesta se revisa con las restricciones obligatorias. Las que no
entran en presupuesto, no caben en los d\'ias disponibles o incumplen una
condici\'on marcada por el usuario no llegan a la salida. Sobre las rutas que
s\'i sobreviven se aplican criterios de comparaci\'on para ordenar alternativas.
Finalmente se presentan dos recomendaciones distintas si existen; en caso
contrario se informa de que el conjunto de soluciones factibles ha quedado
vac\'io.

\section{Formalizaci\'on}
\subsection{Naturaleza del problema y m\'etodo de resoluci\'on escogido}
\subsubsection{Clasificaci\'on del perfil del usuario}
La primera parte del problema encaja con una tarea de \textbf{clasificaci\'on}.
El sistema recibe datos directos del usuario, pero para recomendar necesita
trabajar con un perfil de viaje ya interpretado. El atributo m\'as importante de
esa clasificaci\'on es el motivo del viaje, que puede quedar como descanso,
cultural, diversi\'on, rom\'antico, naturaleza o mixto.

Cuando el usuario ya expresa un motivo claro, esa respuesta se conserva. Si la
respuesta queda abierta, la clasificaci\'on se apoya en el tipo de grupo y en el
evento: una boda o un aniversario orientan el caso hacia el romanticismo, un
fin de curso hacia la diversi\'on, viajar con ni\~nos hacia una propuesta de
descanso y viajar solo sin objetivo expl\'icito hacia una ruta cultural. Tambi\'en
se normalizan otros datos antes de generar rutas, como pasar el presupuesto a
un nivel orientativo y convertir la calidad m\'inima del alojamiento a una
escala num\'erica.

Esta primera tarea reduce la ambig\"uedad de la petici\'on. En vez de generar
planes con respuestas incompletas, el sistema llega a la siguiente fase con un
usuario formalizado mediante atributos que pueden aparecer en condiciones de
reglas.

\subsubsection{Generaci\'on, revisi\'on y comparaci\'on de planes}
La segunda parte del problema encaja con un proceso de
\textbf{propuesta y revisi\'on}. A partir del perfil ya clasificado se proponen
combinaciones de recursos de la base de conocimiento: ciudades, alojamientos,
gu\'ias de visita y conexiones. Cada combinaci\'on intenta formar un itinerario
de dos o tres ciudades.

La revisi\'on se hace comprobando restricciones durante la propia generaci\'on.
No se debe aceptar una propuesta si usa recursos incompatibles con ni\~nos, si
incumple la calidad m\'inima, si necesita un avi\'on que el usuario ha rechazado,
si la distancia no tiene sentido para los d\'ias disponibles o si el precio
calculado supera el presupuesto. Por eso en la implementaci\'on solo se
materializan como \texttt{PlanCandidato} las combinaciones que han sobrevivido
a esos filtros.

El m\'etodo se completa con una \textbf{comparaci\'on heur\'istica}. Las
restricciones determinan qu\'e planes son v\'alidos; la puntuaci\'on permite
ordenar los v\'alidos para preferir rutas con mejor adecuaci\'on base, tren si se
ha pedido, ciudades menos conocidas, calidad por encima del m\'inimo y margen de
presupuesto.

\subsubsection{Encaje entre los subproblemas y el m\'etodo de resoluci\'on}
La descomposici\'on del apartado anterior se puede encajar en el m\'etodo de
resoluci\'on de esta forma:
\begin{itemize}
    \item La adquisici\'on de datos prepara el hecho que representa al usuario.
    \item La deducci\'on del perfil aplica la clasificaci\'on y completa los
    atributos que se usar\'an para decidir.
    \item El filtrado y la generaci\'on de itinerarios forman la parte de
    propuesta y revisi\'on de soluciones.
    \item El c\'alculo de precio revisa la restricci\'on de presupuesto y la
    puntuaci\'on compara los planes factibles.
    \item La selecci\'on final escoge dos recomendaciones distintas a partir del
    conjunto de candidatos ordenados.
    \item La salida muestra la soluci\'on completa o informa de no soluci\'on si
    la revisi\'on ha dejado el conjunto de candidatos vac\'io.
\end{itemize}

Por tanto, no hemos formalizado el sistema como una \'unica regla enorme. La
resoluci\'on se reparte en subtareas que transforman datos de usuario en hechos
de trabajo, hechos candidatos y, finalmente, recomendaciones.

\subsection{Desarrollo de la ontolog\'ia}
\subsubsection{Dominio, cobertura y preguntas que debe responder}
La ontolog\'ia se ha construido para cubrir el conocimiento que hace falta al
recomendar viajes personalizados. El dominio no es el turismo completo, sino
la parte que interviene en nuestros planes: perfil del cliente, destinos,
servicios, restricciones, preferencias, etapas y paquete de viaje resultante.

Para decidir la cobertura partimos de preguntas que el sistema debe poder
responder:
\begin{itemize}
    \item Qu\'e informaci\'on del usuario condiciona la recomendaci\'on.
    \item Qu\'e ciudades, alojamientos, actividades y transportes pueden formar
    una ruta.
    \item Qu\'e recursos pertenecen a cada etapa y qu\'e etapas pertenecen al
    paquete completo.
    \item Qu\'e condiciones debe respetar un viaje y qu\'e preferencias puede
    cumplir sin tratarlas como obligatorias.
    \item Qu\'e precio y qu\'e puntuaci\'on resumen la recomendaci\'on final.
\end{itemize}

A partir de esas preguntas hemos seguido un desarrollo progresivo: primero
identificamos los conceptos del dominio, despu\'es separamos servicios y
soluciones, luego a\~nadimos atributos medibles y por \'ultimo definimos las
relaciones necesarias entre el paquete, sus etapas y los recursos usados.

\subsubsection{T\'erminos relevantes del dominio}
Los t\'erminos que hemos considerado relevantes para pasar de la
conceptualizaci\'on a la ontolog\'ia son:
\begin{itemize}
    \item Cliente o usuario, grupo viajero, edad, presupuesto, duraci\'on,
    objetivo, restricciones y preferencias.
    \item Ciudad o destino, zona, coste diario, seguridad, popularidad y
    etiquetas tur\'isticas.
    \item Alojamiento, categor\'ia, calidad, precio por noche y compatibilidad
    familiar.
    \item Actividad o visita, transporte o conexi\'on y coste asociado.
    \item Etapa, paquete de viaje, precio total y puntuaci\'on.
\end{itemize}

Estos t\'erminos no se han convertido todos en clases. Algunos son atributos,
otros son relaciones y otros se concretan como valores controlados en CLIPS.
La distinci\'on evita que la ontolog\'ia tenga clases innecesarias para cada
valor de una preferencia o de una categor\'ia.

\subsubsection{Clases, jerarqu\'ia y decisiones de modelado}
La jerarqu\'ia parte de la clase general \texttt{Entidad}. Bajo ella situamos
el usuario, la ciudad, la restricci\'on, la preferencia, el paquete de viaje,
la etapa y \texttt{Servicio}. Esta \'ultima clase agrupa los recursos que se
consumen dentro del recorrido. Sus especializaciones son
\texttt{Alojamiento}, \texttt{Transporte} y \texttt{Actividad}.

Esta separaci\'on responde a dos decisiones de modelado. Primero, el cliente y
la soluci\'on no son servicios: el usuario describe el caso y el paquete es el
resultado. Segundo, restricciones y preferencias se mantienen como conceptos
del dominio porque tienen papeles diferentes en la recomendaci\'on, aunque en
la implementaci\'on operativa algunas se representen como atributos del
usuario.

Hemos incluido \texttt{Etapa} porque un paquete de varias ciudades se entiende
mejor como una secuencia de partes: cada etapa visita una ciudad y usa un
alojamiento y unas actividades. En CLIPS esta estructura se compacta dentro
del plan por simplicidad, pero en la ontolog\'ia permite expresar el dominio de
forma m\'as natural.

\subsubsection{Propiedades, atributos y restricciones de valores}
Los atributos de datos recogen la informaci\'on que se necesita medir o mostrar.
En el usuario se representan la edad, el presupuesto en euros y la duraci\'on
m\'axima. En ciudad aparece el coste diario. En alojamiento aparecen el precio
por noche y la calidad. En transporte se guardan coste y horas. Por \'ultimo,
el paquete de viaje tiene precio total y puntuaci\'on.

Adem\'as de esas propiedades de la ontolog\'ia, la formalizaci\'on operativa
necesita valores controlados para poder escribir reglas precisas. Por eso en
CLIPS se restringen, entre otros, los tipos de grupo, los motivos del viaje,
los niveles de calidad, los medios de transporte, la popularidad de una ciudad
y los \'ambitos de una conexi\'on. As\'i evitamos que el razonamiento dependa de
texto libre.

\subsubsection{Relaciones entre conceptos}
Las relaciones de objeto de la ontolog\'ia expresan c\'omo se arma una
recomendaci\'on:
\begin{itemize}
    \item \texttt{viajaA} une un \texttt{PaqueteViaje} con las
    \texttt{Ciudad} que visita.
    \item \texttt{incluyeEtapa} une el paquete con cada \texttt{Etapa}.
    \item \texttt{usaAlojamiento} e \texttt{incluyeActividad} describen los
    recursos de una etapa.
    \item \texttt{usaTransporte} indica los desplazamientos necesarios para el
    paquete completo.
    \item \texttt{cumplePreferencia} y \texttt{respetaRestriccion} distinguen
    los criterios opcionales de las condiciones obligatorias.
\end{itemize}

Con estas relaciones la ontolog\'ia no solo enumera conceptos, sino que permite
leer una recomendaci\'on como un objeto compuesto.

\subsubsection{Instancias necesarias para la base de conocimiento}
Para resolver casos concretos necesitamos instancias de los conceptos
anteriores. En nuestra base de conocimiento se materializan como hechos
iniciales de CLIPS:
\begin{itemize}
    \item Ciudades disponibles con coste, popularidad, compatibilidad familiar,
    puntuaci\'on base y etiquetas.
    \item Alojamientos asociados a una ciudad con calidad y precio por noche.
    \item Actividades y gu\'ias de visita que dan contenido concreto a cada
    etapa.
    \item Conexiones desde el origen y entre ciudades con medio, coste, horas y
    \'ambito.
\end{itemize}

Las instancias se han seleccionado a partir de la informaci\'on de viajes
consultada y se han adaptado a la granularidad del sistema. La ontolog\'ia
define su estructura; la base de hechos aporta los recursos concretos con los
que el motor puede construir rutas.

\subsection{Ontolog\'ia final del sistema}
\subsubsection{Grafo de la jerarqu\'ia de conceptos}
La Figura \ref{fig:jerarquia-ontologia} muestra la jerarqu\'ia final de clases.
La rama de \texttt{Servicio} agrupa los recursos utilizados por el paquete,
mientras que las dem\'as ramas representan el caso, la soluci\'on y los
criterios que se aplican sobre ella.

\begin{figure}[ht]
    \centering
    \begin{tikzpicture}[
        box/.style={draw=UPCBlue, rounded corners=2pt, align=center,
                    font=\footnotesize, inner sep=4pt},
        edge/.style={-{Latex[length=2mm]}, draw=UPCBlue}
    ]
        \node[box] (entidad) at (0,3.0) {Entidad};
        \node[box] (usuario) at (-6.3,1.4) {Usuario};
        \node[box] (ciudad) at (-4.4,1.4) {Ciudad};
        \node[box] (servicio) at (-2.4,1.4) {Servicio};
        \node[box] (restriccion) at (-0.1,1.4) {Restriccion};
        \node[box] (preferencia) at (2.15,1.4) {Preferencia};
        \node[box] (paquete) at (4.65,1.4) {PaqueteViaje};
        \node[box] (etapa) at (6.85,1.4) {Etapa};
        \node[box] (alojamiento) at (-4.5,-0.3) {Alojamiento};
        \node[box] (transporte) at (-2.2,-0.3) {Transporte};
        \node[box] (actividad) at (0.1,-0.3) {Actividad};

        \draw[edge] (entidad) -- (usuario);
        \draw[edge] (entidad) -- (ciudad);
        \draw[edge] (entidad) -- (servicio);
        \draw[edge] (entidad) -- (restriccion);
        \draw[edge] (entidad) -- (preferencia);
        \draw[edge] (entidad) -- (paquete);
        \draw[edge] (entidad) -- (etapa);
        \draw[edge] (servicio) -- (alojamiento);
        \draw[edge] (servicio) -- (transporte);
        \draw[edge] (servicio) -- (actividad);
    \end{tikzpicture}
    \caption{Jerarqu\'ia de conceptos de la ontolog\'ia de viajes.}
    \label{fig:jerarquia-ontologia}
\end{figure}

\subsubsection{Descripci\'on detallada de clases}
Las clases finales tienen el siguiente papel:
\begin{itemize}
    \item \texttt{Entidad}: ra\'iz com\'un de los conceptos principales del
    dominio.
    \item \texttt{Usuario}: perfil del cliente, con datos del grupo, objetivo,
    condiciones y preferencias.
    \item \texttt{Ciudad}: destino candidato que aporta caracter\'isticas
    tur\'isticas y costes.
    \item \texttt{Servicio}: clase general para recursos que se usan dentro de
    una recomendaci\'on.
    \item \texttt{Alojamiento}: servicio de estancia asociado a una ciudad.
    \item \texttt{Transporte}: servicio que permite conectar origen y destinos.
    \item \texttt{Actividad}: visita concreta que puede formar parte de una
    etapa.
    \item \texttt{Restriccion}: condici\'on que un paquete debe respetar para
    ser v\'alido.
    \item \texttt{Preferencia}: criterio opcional que mejora o explica una
    alternativa factible.
    \item \texttt{PaqueteViaje}: soluci\'on recomendada con ruta, precio y
    puntuaci\'on.
    \item \texttt{Etapa}: parte del paquete centrada en una ciudad concreta.
\end{itemize}

\subsubsection{Descripci\'on detallada de atributos y relaciones}
Las propiedades de datos de la ontolog\'ia son \texttt{edad},
\texttt{presupuestoEuros} y \texttt{duracionMaxDias} para el usuario;
\texttt{nombre} para nombrar recursos; \texttt{costeDiario} para ciudad;
\texttt{precioNoche} y \texttt{calidad} para alojamiento;
\texttt{costeTransporte} y \texttt{horasTransporte} para transporte; y
\texttt{precioTotal} y \texttt{puntuacion} para el paquete de viaje.

Las relaciones finales conectan el paquete con la ruta y sus recursos. El
paquete usa \texttt{viajaA}, \texttt{incluyeEtapa} y
\texttt{usaTransporte}; cada etapa usa \texttt{usaAlojamiento} e
\texttt{incluyeActividad}. A ellas se suman \texttt{cumplePreferencia} y
\texttt{respetaRestriccion}, que dejan claro que un paquete se eval\'ua con
dos tipos de criterios.

\subsubsection{Correspondencia entre la ontolog\'ia y la representaci\'on en CLIPS}
La ontolog\'ia y CLIPS representan el mismo dominio con niveles de detalle
distintos. La correspondencia principal es:
\begin{itemize}
    \item \texttt{Usuario}, \texttt{Ciudad}, \texttt{Alojamiento} y
    \texttt{Actividad} se convierten directamente en \texttt{deftemplate} con
    el mismo papel.
    \item \texttt{Transporte} se implementa como \texttt{Conexion}, porque en
    el razonamiento interesa guardar origen, destino, medio, coste, horas y
    \'ambito del tramo.
    \item \texttt{PaqueteViaje} se representa con \texttt{PlanCandidato}; al
    seleccionar una recomendaci\'on aparece adem\'as el hecho
    \texttt{PlanElegido}.
    \item \texttt{Etapa} se aplana dentro del plan con los campos de ciudad,
    d\'ias, alojamiento, visitas y transporte de cada posici\'on de la ruta.
    \item \texttt{Restriccion} y \texttt{Preferencia} se materializan sobre
    todo como atributos del usuario y como condiciones o puntuaciones en las
    reglas.
\end{itemize}

CLIPS a\~nade algunos hechos auxiliares que no son clases del dominio
ontol\'ogico. \texttt{Fase} controla el avance del razonamiento,
\texttt{DecisionExperta} registra deducciones realizadas y
\texttt{GuiaVisitas} empaqueta visitas compatibles para que la recomendaci\'on
final sea concreta.

\subsection{Formalizaci\'on del razonamiento por subproblemas}
\subsubsection{Datos de entrada y hechos de trabajo}
El caso de entrada se formaliza en un hecho \texttt{Usuario}. La entrevista
rellena sus atributos con respuestas validadas: edad, grupo, presencia de
ni\~nos, evento, motivo, presupuesto, duraci\'on, rango de ciudades,
preferencias de transporte, calidad m\'inima, posibles sacrificios, inter\'es
por destinos menos conocidos y ritmo.

El razonamiento combina ese hecho con la base est\'atica de ciudades,
alojamientos, actividades, gu\'ias de visita y conexiones. Como hechos de
trabajo se usan \texttt{Fase} para controlar en qu\'e subproblema estamos,
\texttt{DecisionExperta} para dejar constancia de algunas deducciones,
\texttt{PlanCandidato} para las soluciones construidas y
\texttt{PlanElegido} para las recomendaciones finales.

\subsubsection{Restricciones duras del viaje}
Las restricciones duras se aplican antes de aceptar un plan candidato:
\begin{itemize}
    \item El n\'umero de ciudades debe estar dentro del rango pedido y cada
    etapa debe respetar los d\'ias m\'inimos y m\'aximos manejados por el
    sistema.
    \item El motivo del viaje debe ser compatible con las etiquetas de las
    ciudades y con la gu\'ia de visitas escogida.
    \item Si viajan ni\~nos, ciudades y alojamientos deben estar marcados como
    aptos.
    \item La calidad del alojamiento debe alcanzar el m\'inimo solicitado,
    salvo la relajaci\'on controlada cuando el usuario acepta sacrificar
    calidad.
    \item Si el usuario evita avi\'on, ning\'un tramo de la ruta puede usarlo.
    En viajes cortos tambi\'en se descartan tramos intercontinentales poco
    razonables.
    \item El precio total calculado debe ser menor o igual que el presupuesto
    disponible.
\end{itemize}

Estas condiciones se formalizan como pruebas de compatibilidad y condiciones
num\'ericas durante la generaci\'on de rutas.

\subsubsection{Preferencias y heur\'isticas de puntuaci\'on}
Las preferencias se formalizan como bonificaciones sobre planes que ya son
v\'alidos. De forma resumida, la puntuaci\'on de una ruta se calcula como:
\[
    S = \sum p_{base} + B_{tren} + B_{oculto} + B_{calidad} + B_{precio}.
\]

El bonus de tren solo aparece si el usuario lo prefiere y alg\'un tramo lo usa.
El bonus de ciudades menos conocidas depende de la popularidad de los destinos.
El bonus de calidad premia alojamientos que superan la calidad m\'inima. El
bonus de precio favorece que el plan deje margen respecto al presupuesto. As\'i
el sistema puede comparar soluciones sin transformar cada preferencia en una
prohibici\'on.

\subsubsection{Construcci\'on del precio total y duraci\'on del viaje}
La duraci\'on de un candidato parte de los d\'ias indicados por el usuario y se
reparte entre sus etapas. En rutas de dos ciudades se divide en dos bloques; en
rutas de tres ciudades se distribuye en tres partes. Despu\'es se comprueba que
cada parte respeta los l\'imites de d\'ias por ciudad y el m\'inimo exigido por
la gu\'ia de visitas escogida.

El precio total se formaliza sumando transportes, estancias, coste diario de
cada ciudad y visitas. Para dos ciudades:
\[
    P = T + d_1(a_1 + c_1) + d_2(a_2 + c_2) + v_1 + v_2.
\]

En esa expresi\'on, \(T\) suma los transportes, \(a_i\) es el precio por noche
del alojamiento de la etapa \(i\), \(c_i\) su coste diario y \(v_i\) el coste
de sus visitas. Para tres ciudades se a\~nade el tercer tramo interno, la
tercera estancia y las visitas de la tercera etapa. Esta expresi\'on permite
comparar todos los planes con el presupuesto del mismo modo.

\subsubsection{Elecci\'on de dos recomendaciones diferentes}
La selecci\'on parte de los \texttt{PlanCandidato} ya puntuados. El primer plan
elegido es el candidato con mayor score. Para el segundo se prioriza el mejor
plan cuya ruta no comparte ciudades con el primero. Si no existe una
alternativa sin solape, se busca la mejor ruta distinta, evitando devolver dos
veces el mismo recorrido con una variaci\'on superficial.

Esta formalizaci\'on responde a la necesidad de entregar dos alternativas
reales. La selecci\'on no vuelve a calcular los planes, sino que trabaja sobre
el conjunto de candidatos factibles que ha dejado la revisi\'on anterior.

\subsubsection{Justificaci\'on de preferencias cumplidas y no soluci\'on}
Cada candidato conserva un texto de preferencias cumplidas que resume por qu\'e
destaca una vez superadas las restricciones. En la salida se muestran tambi\'en
precio, duraci\'on, ciudades, d\'ias por ciudad, visitas, alojamientos y
transportes, de modo que la recomendaci\'on no sea una puntuaci\'on opaca.

Si no se ha podido elegir ni siquiera el primer plan, el sistema formaliza el
resultado como \textbf{no soluci\'on}. En ese caso se informa de que las
restricciones introducidas no se pueden satisfacer con la base actual y se
orienta sobre qu\'e condiciones podr\'ian relajarse, como presupuesto, d\'ias,
avi\'on o calidad m\'inima.

\section{Implementaci\'on}

\subsection{Representaci\'on del conocimiento en CLIPS}
\subsubsection{\texttt{deftemplates} y hechos principales}
La representaci\'on del conocimiento empieza con los \texttt{deftemplates}. En
ellos definimos la estructura de los hechos que podr\'an existir en la base de
hechos durante la ejecuci\'on. Los templates principales son:
\begin{itemize}
    \item \texttt{Fase}, que controla en qu\'e parte del razonamiento est\'a el
    sistema: bienvenida, preguntas, deducci\'on, generaci\'on, selecci\'on,
    salida o fin.
    \item \texttt{Usuario}, que almacena los datos del cliente, restricciones y
    preferencias.
    \item \texttt{Ciudad}, \texttt{Alojamiento}, \texttt{Actividad},
    \texttt{GuiaVisitas} y \texttt{Conexion}, que forman la base de
    conocimiento tur\'istica.
    \item \texttt{PlanCandidato}, que representa un itinerario factible
    generado por el sistema.
    \item \texttt{PlanElegido}, que marca cu\'ales de los candidatos se van a
    mostrar como recomendaci\'on final.
    \item \texttt{DecisionExperta}, que guarda algunas deducciones relevantes
    realizadas por las reglas.
\end{itemize}

Hemos usado valores controlados en muchos slots, por ejemplo en el tipo de
grupo, motivo del viaje, calidad de alojamiento, medio de transporte o
popularidad de una ciudad. Esto hace que las reglas puedan comparar valores
concretos y evita errores por texto libre.

\subsubsection{Base de conocimiento de ciudades, alojamientos, visitas y conexiones}
La base de conocimiento se carga mediante varios bloques \texttt{deffacts}. El
primero inicializa el sistema en la fase de bienvenida. Despu\'es aparecen los
hechos que representan los recursos disponibles.

En \texttt{CiudadesDisponibles} se definen destinos con nombre, pa\'is,
continente, zona, coste diario, seguridad, compatibilidad con ni\~nos,
popularidad, puntuaci\'on base y etiquetas tur\'isticas. Estas etiquetas son las
que permiten relacionar una ciudad con motivos como cultura, romanticismo,
descanso, diversi\'on o naturaleza.

En \texttt{AlojamientosDisponibles} cada alojamiento queda asociado a una
ciudad y tiene categor\'ia, calidad num\'erica, precio por noche y aptitud para
ni\~nos. En \texttt{ActividadesDisponibles} se guardan visitas concretas con
tipo, horas, coste, prioridad y compatibilidad familiar. Adem\'as,
\texttt{GuiasDeVisita} agrupa visitas por ciudad y motivo para que la salida
final no sea gen\'erica.

Por \'ultimo, \texttt{ConexionesDisponibles} representa transportes desde el
origen y entre ciudades. Cada conexi\'on incluye origen, destino, medio, \'ambito,
coste y horas. Con esto el sistema puede construir rutas multi-ciudad y no
solo elegir destinos independientes.

\subsubsection{Hechos din\'amicos de usuario, candidatos y planes elegidos}
Durante la ejecuci\'on se crean y modifican hechos din\'amicos. El hecho
\texttt{Usuario} empieza vac\'io en modo interactivo y se va rellenando con las
respuestas de la entrevista. En modo de prueba, los casos autom\'aticos insertan
directamente un \texttt{Usuario} completo y el sistema salta a las deducciones.

Los \texttt{PlanCandidato} se generan solo si una combinaci\'on supera los
filtros de compatibilidad, d\'ias, transporte, alojamiento, gu\'ias y
presupuesto. Cada candidato ya contiene todo lo necesario para imprimirse:
ciudades, d\'ias, alojamientos, categor\'ias, transportes, visitas, precio,
puntuaci\'on y preferencias cumplidas.

Los \texttt{PlanElegido} no duplican toda la informaci\'on del viaje. Solo
guardan el orden de recomendaci\'on y el identificador del candidato elegido.
De esta manera la salida final consulta el candidato original y evita repetir
datos en varios hechos.

\subsection{Funciones auxiliares del sistema}
Las funciones auxiliares se han usado para separar c\'alculos repetidos de las
reglas. Las primeras son \texttt{ask-question} y \texttt{ask-number}, que
validan las respuestas del usuario: una comprueba que la respuesta est\'e entre
los valores permitidos y la otra que el n\'umero est\'e dentro del rango
aceptado.

Despu\'es hay funciones de normalizaci\'on, como \texttt{calidad-a-num} y
\texttt{presupuesto-a-nivel}. La primera transforma calidades como hostal,
est\'andar o lujo a una escala num\'erica. La segunda traduce euros a un nivel
orientativo de presupuesto.

El bloque m\'as importante es el de compatibilidad. Ah\'i se comprueba si el
motivo encaja con cada ciudad, si el tipo de gu\'ia es v\'alido para el destino,
si la calidad del alojamiento cumple el m\'inimo indicado y si los transportes
respetan las restricciones del usuario.

Finalmente est\'an las funciones de puntuaci\'on y explicaci\'on. Calculan los
bonos por tren, ciudades menos conocidas, calidad y precio, y tambi\'en generan
el texto de preferencias cumplidas. Adem\'as, se usan funciones espec\'ificas
para comparar rutas y escoger dos recomendaciones distintas cuando es posible.

\subsection{Reglas y fases del motor de inferencia}
\subsubsection{Bienvenida y entrevista}
La ejecuci\'on empieza en la fase \texttt{bienvenida}. Si no hay ning\'un
\texttt{Usuario}, la regla de inicio imprime el encabezado, crea el usuario
por defecto y cambia la fase a \texttt{preguntas}. A partir de ah\'i se disparan
reglas que preguntan edad, grupo, ni\~nos, evento, motivo, presupuesto,
duraci\'on, n\'umero de ciudades, preferencias de transporte, calidad, posibles
sacrificios, inter\'es por destinos menos conocidos y ritmo.

La entrevista est\'a organizada como reglas peque\~nas. Cada regla solo pregunta
un dato si todav\'ia est\'a en valor desconocido o sin inicializar. Cuando todos
los campos necesarios est\'an completos, una regla final cambia la fase a
\texttt{deduccion}. Para los casos autom\'aticos existe una regla especial que
detecta un usuario ya cargado y evita repetir la entrevista.

\subsubsection{Deducciones expertas}
En la fase \texttt{deduccion} se aplican reglas que completan o ajustan el
perfil del usuario. Primero se calculan valores derivados, como el nivel de
presupuesto y el n\'umero asociado a la calidad m\'inima. Luego se aplican
deducciones de sentido com\'un: boda o aniversario implican un viaje rom\'antico,
fin de curso implica diversi\'on, viajar con ni\~nos orienta hacia descanso,
pareja sin ni\~nos puede ser rom\'antico, amigos j\'ovenes apuntan a diversi\'on y
viajar solo sin motivo se trata como cultural.

Tambi\'en hay ajustes de restricciones. Si el viaje es corto y el usuario hab\'ia
permitido tres ciudades, se limita a dos para evitar una ruta demasiado
apretada. Si el presupuesto es econ\'omico y el usuario acepta sacrificar
calidad, el sistema permite bajar la calidad m\'inima de forma controlada.

\subsubsection{Generaci\'on de candidatos}
La generaci\'on tiene dos reglas principales: una para rutas de dos ciudades y
otra para rutas de tres. Ambas siguen la misma idea. Primero leen el perfil del
usuario, despu\'es buscan ciudades compatibles con el motivo y la presencia de
ni\~nos, conexiones de transporte, alojamientos v\'alidos y gu\'ias de visita
coherentes.

Si todos los recursos encajan, la regla reparte los d\'ias entre las ciudades y
calcula el precio total. El precio incluye transportes, noches de alojamiento,
coste diario de cada ciudad y coste de las visitas. Solo si el precio no supera
el presupuesto se calcula la puntuaci\'on y se afirma un
\texttt{PlanCandidato}. As\'i el sistema no llena la base de hechos con planes
que luego habr\'ia que descartar.

\subsubsection{Selecci\'on de recomendaciones}
Cuando termina la fase de generaci\'on, el sistema pasa a \texttt{seleccion}. La
primera regla escoge como primer plan el candidato con mayor puntuaci\'on. Para
el segundo plan se intenta escoger el mejor candidato sin solape de ciudades
con el primero.

Si no existe una alternativa sin solape completo, hay una regla de respaldo que
busca la mejor ruta distinta. Esta decisi\'on evita imprimir dos veces el mismo
viaje con cambios menores, que ser\'ia una salida pobre para el enunciado.

\subsubsection{Salida explicada y cierre del sistema}
En la fase \texttt{salida} hay dos caminos. Si no se ha elegido ning\'un plan,
se imprime un mensaje de no soluci\'on indicando que las restricciones no pueden
satisfacerse con la base actual y sugiriendo relajar presupuesto, d\'ias, avi\'on
o calidad.

Si hay planes elegidos, el sistema imprime cada recomendaci\'on con puntuaci\'on
interna, precio estimado por persona, duraci\'on, ciudades y d\'ias por ciudad,
visitas, alojamientos, transportes y preferencias cumplidas. Si solo se ha
encontrado un plan factible, se avisa de que convendr\'ia relajar restricciones
para obtener una segunda alternativa. Finalmente se cambia la fase a
\texttt{fin}.

\subsection{Mecanismos de control y modularizaci\'on del razonamiento}
El principal mecanismo de control es el hecho \texttt{Fase}. En vez de dejar
que todas las reglas puedan activarse en cualquier momento, cada grupo de
reglas exige estar en una fase concreta. Esto hace que el razonamiento avance
en orden: entrevista, deducci\'on, generaci\'on, selecci\'on y salida.

Tambi\'en se usan prioridades con \texttt{declare (salience ...)}. Las reglas de
generaci\'on tienen prioridad positiva para crear candidatos antes de pasar a
selecci\'on. Las reglas de transici\'on tienen salience negativa para ejecutarse
cuando ya no quedan reglas m\'as espec\'ificas de esa fase. En selecci\'on, el
primer plan y el segundo sin solape tienen m\'as prioridad que el respaldo con
ruta distinta.

La modularizaci\'on no se ha hecho con varios archivos, sino con secciones
claras dentro del mismo \texttt{viajes.clp}. Esto nos ha parecido suficiente
para el tama\~no de la pr\'actica y facilita cargar el sistema con una sola
orden.

\subsection{Decisiones de implementaci\'on relevantes}
\subsubsection{Restricciones obligatorias frente a preferencias}
Una decisi\'on importante ha sido tratar de forma diferente restricciones y
preferencias. Las restricciones aparecen como condiciones que un candidato debe
cumplir: presupuesto, d\'ias, n\'umero de ciudades, evitar avi\'on, calidad m\'inima
y compatibilidad familiar. Si fallan, el plan ni siquiera se afirma como
\texttt{PlanCandidato}.

Las preferencias se aplican despu\'es mediante puntuaci\'on. Por ejemplo, usar
tren, incluir ciudades menos masificadas, superar la calidad m\'inima o dejar
margen de presupuesto suma puntos, pero no invalida por s\'i solo un plan que
cumple las restricciones. Esta separaci\'on hace que el sistema sea menos
fr\'agil y permita comparar alternativas razonables.

\subsubsection{C\'alculo num\'erico de precios}
Hemos usado precios num\'ericos porque el enunciado pide mostrar una
recomendaci\'on concreta y porque el presupuesto es una restricci\'on dura. Cada
candidato suma el coste de los transportes, el coste de alojamiento por noche,
el coste diario de cada ciudad y el coste de las gu\'ias de visita.

El c\'alculo se hace dentro de las reglas de generaci\'on, justo antes de afirmar
el plan. Esto evita generar planes que ya se sabe que superan el presupuesto.
Adem\'as, el margen restante se usa como parte de la puntuaci\'on con
\texttt{bonus-precio}, de manera que dos planes v\'alidos pueden compararse
teniendo en cuenta cu\'al aprovecha mejor el presupuesto.

\subsubsection{Control de rutas distintas y caso imposible}
Para cumplir la idea de dos recomendaciones distintas, el identificador del
plan incluye ciudades, alojamientos, transportes y gu\'ias de visita. Adem\'as,
la selecci\'on del segundo plan usa funciones espec\'ificas para comprobar si dos
rutas no comparten ciudades o, como m\'inimo, si no son exactamente la misma.

El caso imposible se controla en la salida. Si despu\'es de generar y seleccionar
no existe ning\'un \texttt{PlanElegido} de orden 1, se dispara la regla de no
soluci\'on. Esto es mejor que forzar un viaje incorrecto, porque mantiene la
coherencia entre restricciones y recomendaci\'on.

\subsection{Desarrollo incremental y prototipos}
\subsubsection{Prototipo inicial}
El primer objetivo fue tener un prototipo peque\~no pero completo: preguntar
datos del usuario, deducir un perfil b\'asico y producir alguna recomendaci\'on.
En esta fase lo importante no era tener toda la base de conocimiento, sino
comprobar que CLIPS pod\'ia encadenar preguntas, hechos de usuario y reglas de
recomendaci\'on.

Esa idea se conserva en la versi\'on final con la fase de entrevista, las reglas
de deducci\'on y el control por \texttt{Fase}. Aunque el sistema final es m\'as
grande, sigue teniendo el mismo esqueleto que un prototipo funcional.

\subsubsection{Extensiones y revisiones intermedias}
Despu\'es ampliamos el prototipo para cubrir mejor lo que ped\'ia el enunciado.
Las extensiones principales fueron:
\begin{itemize}
    \item pasar de destinos sueltos a itinerarios de dos o tres ciudades;
    \item a\~nadir transporte entre etapas y vuelta al origen;
    \item incorporar alojamientos, visitas concretas y gu\'ias de visita;
    \item calcular precios num\'ericos;
    \item separar restricciones de preferencias;
    \item seleccionar dos planes diferentes;
    \item incluir casos de no soluci\'on y juegos de prueba autom\'aticos.
\end{itemize}

Estas revisiones hicieron que el sistema dejara de ser una recomendaci\'on
simple y pasara a generar paquetes de viaje completos.

\subsubsection{Versi\'on final}
La versi\'on final integra todos los bloques en un \'unico flujo de razonamiento.
El usuario puede usar el sistema de forma interactiva o ejecutar directamente
un caso de prueba. En ambos modos se llega a la misma secuencia: deducci\'on del
perfil, generaci\'on de candidatos, selecci\'on de recomendaciones y salida
explicada.

El resultado final cumple los puntos principales que necesit\'abamos: rutas de
varias ciudades, transportes entre etapas, precio total, visitas concretas,
alojamientos, dos alternativas diferentes cuando sea posible y mensaje de no
soluci\'on cuando las restricciones son incompatibles.

\subsubsection{Planificaci\'on y organizaci\'on del trabajo del grupo}
Nos hemos organizado dividiendo el trabajo por partes del sistema. Primero se
defini\'o el dominio y la ontolog\'ia, despu\'es se implement\'o la estructura en
CLIPS y finalmente se completaron reglas, base de conocimiento y pruebas. Esta
separaci\'on nos permiti\'o avanzar sin mezclar todo a la vez.

\subsection{Instrucciones de ejecuci\'on}
Para ejecutar el sistema en modo interactivo desde CLIPS:
\begin{verbatim}
(clear)
(load "viajes.clp")
(reset)
(run)
\end{verbatim}

Tambi\'en se pueden ejecutar los juegos de prueba definidos al final del archivo:
\begin{verbatim}
(clear)
(load "viajes.clp")
(caso-familia-alto)
(caso-estudiante-tren)
(caso-pareja-romantica)
(caso-amigos-diversion)
(caso-naturaleza-menos-conocido)
(caso-imposible-sin-avion-asia)
\end{verbatim}

Estas pruebas no son una ejecuci\'on interactiva en la que vamos respondiendo
por teclado. Son funciones ya preparadas al final de \texttt{viajes.clp} que
hacen \texttt{reset}, insertan un usuario con las respuestas ya cargadas y
despu\'es ejecutan \texttt{run}. A efectos del razonamiento act\'uan igual que si
ese usuario hubiera contestado la entrevista completa, porque se activa el
mismo motor de inferencia, las mismas reglas y la misma salida final.

\section{Juegos de prueba y validaci\'on}
\subsection{Criterios de selecci\'on de los juegos de prueba}
Los juegos de prueba se han escogido para comprobar que el sistema no funciona
solo con un caso favorable, sino con perfiles distintos y restricciones
diferentes. Hemos intentado cubrir viajes familiares, viajes con presupuesto
ajustado, deducciones de perfil, preferencias de transporte, rutas de ocio,
viajes de naturaleza y un caso imposible.

Cada prueba se ejecuta con una funci\'on definida al final de
\texttt{viajes.clp}. No son casos interactivos escritos a mano durante la
ejecuci\'on, sino funciones que simulan esa interacci\'on: hacen \texttt{reset},
insertan un usuario ya completo y ejecutan \texttt{run}. De esta forma se prueba
exactamente el mismo motor de inferencia que en modo interactivo, pero sin
tener que contestar manualmente a todas las preguntas cada vez.

\subsection{Cobertura de casos comunes, restricciones y excepciones}
El conjunto de pruebas cubre los siguientes aspectos del sistema:
\begin{itemize}
    \item Recomendaci\'on para una familia con ni\~nos y calidad alta.
    \item Restricci\'on dura de evitar avi\'on y preferencia por tren.
    \item Deducci\'on de viaje rom\'antico a partir de un aniversario.
    \item Deducci\'on de diversi\'on a partir de un fin de curso.
    \item Preferencia por naturaleza y lugares menos conocidos.
    \item Detecci\'on del caso sin recomendaci\'on factible.
\end{itemize}

La salida completa de las pruebas se conserva en
\texttt{SALIDAS\_PRUEBAS.txt}. En este apartado recogemos los datos principales
de esa salida para justificar el comportamiento del sistema.

\subsection{Caso 1: familia con ni\~nos y presupuesto alto}
\subsubsection{Escenario, objetivo y entrada}
Este caso representa una familia con ni\~nos, presupuesto alto y preferencia por
un viaje de descanso. El usuario pide entre dos y tres ciudades, ocho d\'ias de
duraci\'on, calidad m\'inima superior y no impone restricciones sobre el avi\'on.

El objetivo de la prueba es comprobar que el sistema respeta la compatibilidad
familiar, mantiene alojamientos de calidad alta y genera un itinerario
multi-ciudad dentro del presupuesto.

Comando ejecutado:
\begin{verbatim}
(caso-familia-alto)
\end{verbatim}

\subsubsection{Resultado esperado}
Esperamos obtener dos planes factibles de descanso, con ciudades y alojamientos
aptos para ni\~nos, precio inferior a 4500 euros por persona y visitas
coherentes con un viaje familiar o relajado.

\subsubsection{Salida original de CLIPS}
La salida genera dos recomendaciones. Los valores principales fueron:
\begin{verbatim}
Plan 1: Punta Cana, Cancun, Costa Rica
Precio: 3488 euros
Duracion: 8 dias
Puntuacion: 237
Preferencias: respeta la calidad minima de alojamiento

Plan 2: Cancun, Costa Rica
Precio: 3222 euros
Duracion: 8 dias
Puntuacion: 169
Preferencias: respeta la calidad minima de alojamiento
\end{verbatim}

\subsubsection{Explicaci\'on del resultado obtenido}
El resultado es correcto porque ambos planes quedan por debajo del presupuesto
y usan alojamientos familiares de categor\'ia alta, como resorts o hoteles
superiores. Adem\'as, el sistema propone destinos de playa, descanso y
naturaleza, que encajan con el motivo del viaje y con el perfil familiar. La
primera recomendaci\'on tiene m\'as puntuaci\'on porque incluye tres ciudades y
combina m\'as recursos compatibles.

\subsection{Caso 2: estudiante que prioriza tren y evita avi\'on}
\subsubsection{Escenario, objetivo y entrada}
Este caso representa un usuario joven que viaja solo, con presupuesto ajustado,
motivo cultural, preferencia por tren y restricci\'on expl\'icita de evitar
avi\'on. La duraci\'on es de seis d\'ias y el viaje debe tener dos ciudades.

El objetivo es comprobar que evitar avi\'on act\'ua como restricci\'on dura y que
la preferencia por tren se refleja en la puntuaci\'on y en la explicaci\'on final.

Comando ejecutado:
\begin{verbatim}
(caso-estudiante-tren)
\end{verbatim}

\subsubsection{Resultado esperado}
Esperamos que no aparezca ning\'un tramo en avi\'on, que los planes usen
transporte ferroviario y que el presupuesto de 1250 euros sea suficiente para
generar rutas culturales con alojamientos econ\'omicos.

\subsubsection{Salida original de CLIPS}
La salida genera dos planes dentro del presupuesto:
\begin{verbatim}
Plan 1: Lisboa, Sevilla
Precio: 914 euros
Duracion: 6 dias
Puntuacion: 189
Transportes: tren, tren, tren

Plan 2: Paris, Amsterdam
Precio: 1075 euros
Duracion: 6 dias
Puntuacion: 187
Transportes: tren, tren, tren
\end{verbatim}

\subsubsection{Explicaci\'on del resultado obtenido}
El sistema cumple la restricci\'on principal porque todos los transportes son en
tren. Tambi\'en respeta el presupuesto y propone visitas culturales en todas las
ciudades. El primer plan queda ligeramente por encima en puntuaci\'on porque,
adem\'as de usar tren, incluye ciudades menos masificadas y deja m\'as margen de
presupuesto.

\subsection{Caso 3: pareja rom\'antica por aniversario}
\subsubsection{Escenario, objetivo y entrada}
Este caso representa una pareja sin ni\~nos que viaja por aniversario. El motivo
se deja como desconocido para comprobar si el sistema deduce correctamente que
el perfil debe ser rom\'antico. El presupuesto es de 2200 euros, la duraci\'on es
de siete d\'ias y se permiten entre dos y tres ciudades.

Comando ejecutado:
\begin{verbatim}
(caso-pareja-romantica)
\end{verbatim}

\subsubsection{Resultado esperado}
Esperamos que la fase de deducci\'on interprete el aniversario como viaje
rom\'antico y que las rutas propuestas incluyan ciudades y gu\'ias compatibles
con ese motivo. Tambi\'en esperamos que se respete el presupuesto y que la
preferencia por tren se tenga en cuenta cuando sea posible.

\subsubsection{Salida original de CLIPS}
La salida confirma dos recomendaciones rom\'anticas:
\begin{verbatim}
Plan 1: Paris, Roma, Viena
Precio: 1899 euros
Duracion: 7 dias
Puntuacion: 288
Visitas: Torre Eiffel, Trastevere, Belvedere

Plan 2: Lisboa, Sevilla
Precio: 990 euros
Duracion: 7 dias
Puntuacion: 196
Visitas: Alfama, Santa Cruz, Guadalquivir
\end{verbatim}

\subsubsection{Explicaci\'on del resultado obtenido}
El comportamiento es el esperado porque el sistema deduce el motivo rom\'antico
a partir del aniversario y selecciona gu\'ias coherentes con ese perfil. El
primer plan obtiene una puntuaci\'on mucho mayor porque combina tres ciudades
con alta puntuaci\'on base y varias etapas compatibles. El segundo plan ofrece
una alternativa m\'as cercana y barata, tambi\'en dentro del presupuesto.

\subsection{Caso 4: grupo de amigos en viaje de fin de curso}
\subsubsection{Escenario, objetivo y entrada}
Este caso representa un grupo de amigos joven con evento de fin de curso. El
motivo se deja como desconocido para que el sistema deduzca diversi\'on. El viaje
tiene seis d\'ias, dos ciudades, calidad m\'inima est\'andar y presupuesto de 1700
euros.

Comando ejecutado:
\begin{verbatim}
(caso-amigos-diversion)
\end{verbatim}

\subsubsection{Resultado esperado}
Esperamos que el sistema deduzca un perfil de ocio y diversi\'on, proponga
ciudades con actividades nocturnas o de ambiente joven y mantenga el precio por
debajo del presupuesto.

\subsubsection{Salida original de CLIPS}
La salida genera dos alternativas de ocio:
\begin{verbatim}
Plan 1: Lisboa, Ibiza
Precio: 1280 euros
Duracion: 6 dias
Puntuacion: 159
Visitas: Barrio Alto, LX Factory, Cala Comte, Dalt Vila

Plan 2: Lisboa, Sevilla
Precio: 909 euros
Duracion: 6 dias
Puntuacion: 157
Visitas: Barrio Alto, LX Factory, Triana, Alameda
\end{verbatim}

\subsubsection{Explicaci\'on del resultado obtenido}
La prueba valida la deducci\'on de diversi\'on a partir del fin de curso. Los dos
planes incluyen actividades de ocio y ciudades con etiquetas compatibles. En
este caso las dos alternativas comparten Lisboa, pero no son el mismo recorrido:
el sistema recurre a una ruta distinta cuando no encuentra una alternativa
mejor sin solape completo.

\subsection{Caso 5: naturaleza y lugares menos conocidos}
\subsubsection{Escenario, objetivo y entrada}
Este caso representa una pareja que busca naturaleza, acepta avi\'on y prefiere
lugares menos conocidos. El presupuesto es de 3500 euros, la duraci\'on de ocho
d\'ias y se permiten dos o tres ciudades.

Comando ejecutado:
\begin{verbatim}
(caso-naturaleza-menos-conocido)
\end{verbatim}

\subsubsection{Resultado esperado}
Esperamos planes con destinos de naturaleza, visitas relacionadas con ese
motivo y bonificaci\'on por incluir lugares menos masificados. El precio total
debe mantenerse dentro del presupuesto.

\subsubsection{Salida original de CLIPS}
La salida obtiene dos rutas compatibles:
\begin{verbatim}
Plan 1: Kyoto, Bali
Precio: 3442 euros
Duracion: 8 dias
Puntuacion: 178
Visitas: Arashiyama, jardines de templos, Ubud, arrozales

Plan 2: Cancun, Costa Rica
Precio: 2938 euros
Duracion: 8 dias
Puntuacion: 175
Visitas: Playa Delfines, Volcan Arenal, Monteverde
\end{verbatim}

\subsubsection{Explicaci\'on del resultado obtenido}
El resultado es coherente porque las ciudades propuestas tienen etiquetas o
gu\'ias compatibles con naturaleza y descanso. Adem\'as, ambas salidas indican
que se cumple la preferencia por ciudades menos masificadas y la calidad
m\'inima del alojamiento. La primera recomendaci\'on queda m\'as cerca del l\'imite
de presupuesto, pero sigue siendo factible.

\subsection{Caso 6: restricciones incompatibles y no soluci\'on}
\subsubsection{Escenario, objetivo y entrada}
Este caso fuerza una situaci\'on complicada: pareja en viaje de boda, motivo
rom\'antico, solo cuatro d\'ias, presupuesto de 1200 euros, preferencia por tren,
restricci\'on de evitar avi\'on, calidad m\'inima de lujo y sin aceptar sacrificar
calidad ni duraci\'on.

Comando ejecutado:
\begin{verbatim}
(caso-imposible-sin-avion-asia)
\end{verbatim}

\subsubsection{Resultado esperado}
Esperamos que no se genere ninguna recomendaci\'on factible. Las condiciones son
demasiado restrictivas para la base de conocimiento: pocos d\'ias, presupuesto
bajo, calidad de lujo, dos o tres ciudades y prohibici\'on de avi\'on.

\subsubsection{Salida original de CLIPS}
La salida confirma el caso de no soluci\'on:
\begin{verbatim}
NO SE HA ENCONTRADO NINGUNA RECOMENDACION FACTIBLE
Las restricciones introducidas no pueden satisfacerse
con la base de conocimiento actual.
Sugerencias: aumentar presupuesto/dias, permitir avion
o bajar la calidad minima.
\end{verbatim}

\subsubsection{Explicaci\'on del resultado obtenido}
La prueba valida que el sistema no fuerza una recomendaci\'on incorrecta. Si no
existe ning\'un candidato que cumpla todas las restricciones, se informa de
forma clara y se indican posibles condiciones que el usuario podr\'ia relajar.

\subsection{Verificaci\'on del cumplimiento del enunciado}
Con los seis casos se comprueba que el sistema cumple los requisitos principales
del enunciado. El sistema pregunta o recibe datos del usuario, deduce perfiles
cuando falta informaci\'on, genera viajes de varias ciudades, incluye
alojamientos, visitas y transportes, calcula precio total, devuelve dos planes
cuando existen alternativas y contempla el caso sin soluci\'on.

Tambi\'en se verifica la diferencia entre restricciones y preferencias. Evitar
avi\'on, presupuesto, d\'ias, calidad m\'inima y compatibilidad familiar act\'uan
como condiciones obligatorias. En cambio, usar tren, preferir ciudades menos
conocidas o mejorar calidad se usan para puntuar y explicar planes factibles.

\subsection{An\'alisis global de los resultados de prueba}
Los resultados son consistentes con el funcionamiento esperado. En los casos
factibles se obtienen rutas completas, siempre con precio, duraci\'on, ciudades,
visitas, alojamientos, transportes y preferencias cumplidas. En el caso
imposible no se imprime una ruta artificial, sino un mensaje expl\'icito de no
recomendaci\'on.

La prueba de amigos muestra una limitaci\'on razonable del sistema: cuando no
hay suficiente variedad sin solape completo, la segunda alternativa puede
compartir alguna ciudad con la primera. Aun as\'i, no se devuelve el mismo plan,
sino una ruta distinta. Esto mantiene la salida dentro de la base disponible y
evita inventar recursos que no existen.

\section{Conclusiones}
\subsection{Cumplimiento de objetivos}
El sistema final cumple los objetivos principales de la pr\'actica. Hemos
construido un Sistema Basado en el Conocimiento capaz de recomendar viajes
personalizados a partir de datos del usuario, restricciones y preferencias. La
soluci\'on no se limita a elegir una ciudad, sino que genera itinerarios con dos
o tres destinos, transportes, alojamientos, visitas, duraci\'on, precio y una
explicaci\'on de las preferencias cumplidas.

Tambi\'en se cumple el requisito de ofrecer dos alternativas cuando la base de
conocimiento lo permite. Adem\'as, el sistema contempla el caso en que no existe
ninguna combinaci\'on factible y lo comunica de forma expl\'icita.

\subsection{Fortalezas del sistema construido}
Una de las fortalezas del sistema es que separa bien el conocimiento del dominio
del razonamiento. Las ciudades, alojamientos, visitas y conexiones est\'an
representadas como hechos, mientras que las decisiones de experto aparecen en
reglas y funciones de compatibilidad.

Otra fortaleza es la explicabilidad. La salida no muestra solo un nombre de
destino, sino el precio, los d\'ias, los transportes, las visitas, los
alojamientos y las preferencias satisfechas. Esto hace que el usuario pueda
entender por qu\'e se le recomienda cada plan.

Tambi\'en consideramos positiva la modularizaci\'on por fases. El hecho
\texttt{Fase} permite que la entrevista, la deducci\'on, la generaci\'on, la
selecci\'on y la salida se ejecuten en orden y sean m\'as f\'aciles de revisar.

\subsection{Limitaciones actuales y posibles mejoras}
La principal limitaci\'on es que la base de conocimiento es cerrada. El sistema
no consulta disponibilidad real de vuelos, hoteles o actividades, y los precios
son estimaciones usadas para razonar dentro de la pr\'actica. En una versi\'on
m\'as completa se podr\'ian conectar fuentes externas o actualizar precios seg\'un
temporada.

Otra limitaci\'on es que el origen se modela de forma gen\'erica. Esto simplifica
la pr\'actica, pero en un sistema real el origen del usuario afectar\'ia mucho al
precio y al transporte. Tambi\'en podr\'ian incorporarse fechas, clima, visados,
horarios y preferencias m\'as finas sobre ritmo o tipo de alojamiento.

Como mejora del razonamiento, se podr\'ia enriquecer la explicaci\'on final con
las decisiones expertas deducidas durante la ejecuci\'on, por ejemplo indicando
que un aniversario se ha interpretado como viaje rom\'antico o que un viaje
corto ha limitado el n\'umero de ciudades.

\subsection{Aprendizajes sobre ontolog\'ias, reglas y Sistemas Basados en el Conocimiento}
La pr\'actica nos ha servido para ver que construir un SBC no consiste solo en
programar reglas sueltas. Primero hay que entender el dominio, decidir qu\'e
conceptos importan, separar restricciones de preferencias y formalizar una
ontolog\'ia que despu\'es pueda trasladarse a hechos y reglas.

Tambi\'en hemos visto la importancia de probar casos variados. Un sistema puede
parecer correcto con un perfil sencillo, pero fallar cuando se combinan
restricciones fuertes. Por eso los juegos de prueba han sido \'utiles para
comprobar tanto recomendaciones normales como el caso sin soluci\'on.

Finalmente, hemos aprendido que CLIPS encaja bien cuando el problema se puede
expresar mediante hechos, condiciones y reglas de inferencia. En nuestro caso,
el enfoque nos ha permitido representar conocimiento expl\'icito y justificar
las recomendaciones de una forma bastante clara.

\appendix
\section{Salidas literales completas de los juegos de prueba}
Este ap\'endice recoge las salidas completas guardadas en
\texttt{SALIDAS\_PRUEBAS.txt}. En el apartado 6 hemos usado un resumen de cada
caso para no cortar el razonamiento principal de la memoria; aqu\'i dejamos la
evidencia literal que permite comprobar que las pruebas se han ejecutado sobre
el sistema final. Todas estas salidas vienen de funciones de prueba definidas
al final de \texttt{viajes.clp}. Esas funciones precargan el usuario y lanzan
\texttt{run}, por lo que no requieren entrevista manual, pero recorren el mismo
razonamiento que una ejecuci\'on interactiva.

\subsection{Preparaci\'on de la ejecuci\'on}
\lstinputlisting[style=salidaclip,firstline=1,lastline=7]{../SALIDAS_PRUEBAS.txt}

\subsection{Caso familia}
\lstinputlisting[style=salidaclip,firstline=8,lastline=57]{../SALIDAS_PRUEBAS.txt}

\subsection{Caso estudiante}
\lstinputlisting[style=salidaclip,firstline=58,lastline=105]{../SALIDAS_PRUEBAS.txt}

\subsection{Caso pareja}
\lstinputlisting[style=salidaclip,firstline=106,lastline=156]{../SALIDAS_PRUEBAS.txt}

\subsection{Caso amigos}
\lstinputlisting[style=salidaclip,firstline=157,lastline=204]{../SALIDAS_PRUEBAS.txt}

\subsection{Caso naturaleza}
\lstinputlisting[style=salidaclip,firstline=205,lastline=252]{../SALIDAS_PRUEBAS.txt}

\subsection{Caso sin soluci\'on}
\lstinputlisting[style=salidaclip,firstline=253,lastline=267]{../SALIDAS_PRUEBAS.txt}

\section{Documentaci\'on de la ontolog\'ia generada con Prot\'eg\'e}
\subsection{Archivo y alcance de la ontolog\'ia}
La ontolog\'ia se entrega en el archivo \texttt{ontologia\_viajes.ttl}. Est\'a en
formato Turtle y se puede abrir directamente con Prot\'eg\'e. Su objetivo es
documentar el vocabulario conceptual del sistema: las clases principales del
dominio, las propiedades que relacionan los conceptos y los atributos m\'as
importantes que despu\'es se usan en CLIPS.

Hemos separado la ontolog\'ia del motor de inferencia porque cada parte tiene
una funci\'on distinta. La ontolog\'ia sirve para dejar claro qu\'e conceptos
existen y c\'omo se relacionan. El archivo \texttt{viajes.clp}, en cambio,
materializa esa conceptualizaci\'on en hechos y reglas ejecutables.

\subsection{Grafo de clases y relaciones}
La jerarqu\'ia de clases de la ontolog\'ia parte de una clase ra\'iz
\texttt{Entidad}. De ella cuelgan los conceptos principales del dominio:
\texttt{Usuario}, \texttt{Ciudad}, \texttt{Servicio}, \texttt{Restriccion},
\texttt{Preferencia}, \texttt{PaqueteViaje} y \texttt{Etapa}. Dentro de
\texttt{Servicio} hemos agrupado los recursos concretos que puede utilizar un
viaje: \texttt{Alojamiento}, \texttt{Transporte} y \texttt{Actividad}.

\begin{center}
\begin{tabular}{p{0.28\textwidth} p{0.28\textwidth} p{0.34\textwidth}}
Origen & Relaci\'on & Destino \\
\hline
\texttt{PaqueteViaje} & \texttt{viajaA} & \texttt{Ciudad} \\
\texttt{PaqueteViaje} & \texttt{incluyeEtapa} & \texttt{Etapa} \\
\texttt{PaqueteViaje} & \texttt{usaTransporte} & \texttt{Transporte} \\
\texttt{Etapa} & \texttt{usaAlojamiento} & \texttt{Alojamiento} \\
\texttt{Etapa} & \texttt{incluyeActividad} & \texttt{Actividad} \\
\texttt{PaqueteViaje} & \texttt{cumplePreferencia} & \texttt{Preferencia} \\
\texttt{PaqueteViaje} & \texttt{respetaRestriccion} & \texttt{Restriccion} \\
\end{tabular}
\end{center}

Estas relaciones reflejan la estructura real de la recomendaci\'on. Un paquete
de viaje no contiene solo un destino, sino varias etapas, transportes,
alojamientos, actividades y criterios de evaluaci\'on.

\subsection{Fichas de clases, propiedades e instancias}
\subsubsection{Clases principales}
\begin{itemize}
    \item \texttt{Usuario}: representa el perfil del cliente. Incluye edad,
    grupo, presencia de ni\~nos, presupuesto, duraci\'on, motivo, calidad m\'inima
    y preferencias.
    \item \texttt{Ciudad}: representa cada destino candidato. Guarda nombre,
    pa\'is, continente, zona, coste diario, seguridad, compatibilidad familiar,
    popularidad y etiquetas tur\'isticas.
    \item \texttt{Servicio}: clase general para recursos usados por el viaje.
    Sus subclases son \texttt{Alojamiento}, \texttt{Transporte} y
    \texttt{Actividad}.
    \item \texttt{Alojamiento}: recurso de estancia asociado a una ciudad, con
    categor\'ia, calidad, precio por noche y compatibilidad con ni\~nos.
    \item \texttt{Transporte}: conexi\'on utilizada para llegar a los destinos,
    moverse entre ciudades o volver al origen.
    \item \texttt{Actividad}: visita o experiencia concreta de una ciudad.
    \item \texttt{Restriccion}: condici\'on obligatoria que un plan debe cumplir.
    En el sistema aparecen como presupuesto, duraci\'on, evitar avi\'on, n\'umero
    de ciudades o calidad m\'inima.
    \item \texttt{Preferencia}: criterio usado para comparar soluciones
    factibles, por ejemplo usar tren, escoger lugares menos conocidos o mejorar
    la calidad.
    \item \texttt{PaqueteViaje}: soluci\'on recomendada, con ciudades, etapas,
    precio total y puntuaci\'on.
    \item \texttt{Etapa}: parte de un paquete centrada en una ciudad concreta.
\end{itemize}

\subsubsection{Propiedades de objeto}
Las propiedades de objeto indican c\'omo se conectan las clases. La ontolog\'ia
define \texttt{viajaA}, \texttt{incluyeEtapa}, \texttt{usaAlojamiento},
\texttt{incluyeActividad}, \texttt{usaTransporte}, \texttt{cumplePreferencia}
y \texttt{respetaRestriccion}. Con ellas se puede expresar que una
recomendaci\'on visita ciudades, se compone de etapas, usa recursos y se eval\'ua
con restricciones y preferencias.

\subsubsection{Propiedades de datos}
Las propiedades de datos recogen los atributos num\'ericos o textuales
necesarios para razonar. En \texttt{Usuario} aparecen \texttt{edad},
\texttt{presupuestoEuros} y \texttt{duracionMaxDias}. En \texttt{Ciudad}
aparece \texttt{costeDiario}. En \texttt{Alojamiento} se usan
\texttt{precioNoche} y \texttt{calidad}. En \texttt{Transporte} se usan
\texttt{costeTransporte} y \texttt{horasTransporte}. Finalmente,
\texttt{PaqueteViaje} tiene \texttt{precioTotal} y \texttt{puntuacion}.

\subsubsection{Instancias representadas en la base de conocimiento}
El archivo \texttt{ontologia\_viajes.ttl} documenta clases y propiedades. Las
instancias operativas se han representado en \texttt{viajes.clp}, porque son
las que necesita el motor de inferencia durante la ejecuci\'on. La base de
conocimiento contiene:
\begin{itemize}
    \item 16 ciudades en \texttt{CiudadesDisponibles}.
    \item 34 alojamientos en \texttt{AlojamientosDisponibles}.
    \item 32 actividades en \texttt{ActividadesDisponibles}.
    \item 36 gu\'ias de visita en \texttt{GuiasDeVisita}.
    \item 62 conexiones de transporte en \texttt{ConexionesDisponibles}.
\end{itemize}

Esta decisi\'on evita duplicar toda la base de datos en dos formatos distintos.
La ontolog\'ia queda como documentaci\'on conceptual y CLIPS queda como
representaci\'on ejecutable del conocimiento.

\end{document}
